\scanpage{11}\mspage{10}%
ci inspirent.

\uline{Axiomes}\quad On peut prendre pour ens de base \unsure{ds|d'} une axiomatique des ateliers, soit $\unsure{\mathcal{F}|\mathfrak{F}}$, soit $\unsure{\mathcal{M}|M}$, soit $\mathcal{L}$. On trouve des données les plus simples sur $\mathcal{F}$, les \deleted{plus} \unsure{sophistiqués|sophistiquées} sur $\mathcal{L}$. Par contre, les axiomes essentiels sont les moins transparents quand on \unsure{reste|entre} plus ou moins dans $\mathcal{F}$, le plus transparent quand on travaille dans $\mathcal{L}$ (ou peut-être pareil pour $\mathcal{M}$ ?).

\uline{Sur $\mathcal{F}$}, on peut prendre comme structure de base, au choix
\[
\left\{\begin{aligned}
 &\text{a)}\ \leq,\ \ll\\
 &\text{b)}\ \leq,\ \preccurlyeq \quad \longleftarrow\\
 &\text{c)}\ \ll,\ \deleted{\unsure{F|\mathcal F}}\quad
 \underbrace{\text{et}\quad\mathcal{F}\longrightarrow\mathrm{P}(\mathcal{F})}_{
 \begin{smallmatrix}
 \text{i.e. la \uline{relation}}\\[-1pt]
 \text{\uline{d'incidence}}\\
 \text{``$F$ est une strate de $G$''}\\
 \text{(i.e. $F\in\mathcal{M}$ \uline{et} $F\leq G$)}
 \end{smallmatrix}}
\end{aligned}\right.
\quad
\begin{aligned}
 &\text{(notée }F\longmapsto\widetilde F,\ \text{le}\\
 &\text{déployement de }F\\
 &\text{vu comme partie de }\mathcal{F}\text{)}\\
 &\text{\deleted{i.e. la \deleted{\unsure{description|désignation}} \inserted{relation d'inc.}}}
\end{aligned}
\tag*{(1.16)}
\]
% The leftward pointer at b) belongs to this later dated note.
% slanted
\margin{(20. juin) Non, par mon nouveau point de vue ça n'implique plus $\ll$}
% In the canceled note at right of c), the lower word was locally struck,
% an upper replacement was entered, and the entire note was then struck.

\deleted{Je} Avant de préciser \unsure{ces points|ce point}, je vais grouper dans un diagramme les six relations fondamentales dans $\mathcal{F}$, en écrivant
\[
F\triangleleft G\quad
\underset{\unsure{\mathrm{par}}}{\overset{\mathrm{dfn}}{\Longleftrightarrow}}\quad
\text{``$F$ incident à $G$''}\quad\text{i.e.}\quad F\leq G\ \text{et}\ F\in\mathcal{M}
\tag*{(1.17)}
\]
